% ---------------------------------------------------------------
% AI-ENG-110 - Final Project Report - Topic 4
% Compile with: pdflatex final_report.tex  (run twice for the ToC)
% ---------------------------------------------------------------
\documentclass[11pt,a4paper]{article}

\usepackage[T1]{fontenc}
\usepackage[utf8]{inputenc}
\usepackage[margin=2.5cm]{geometry}
\usepackage{lmodern}
\usepackage{microtype}
\usepackage{booktabs}
\usepackage{longtable}
\usepackage{array}
\usepackage{enumitem}
\usepackage{fancyhdr}
\usepackage{listings}
\usepackage{xcolor}
\usepackage{titlesec}
\usepackage[hidelinks]{hyperref}

% ---------- header / footer ----------
\pagestyle{fancy}
\fancyhf{}
\fancyhead[L]{\small AI-ENG-110 -- Final Report -- Topic 4}
\fancyhead[R]{\small Spring 2026}
\fancyfoot[L]{\small AI Academy, National AI Center}
\fancyfoot[R]{\small Page \thepage\ of \pageref{LastPage}}
\renewcommand{\headrulewidth}{0.4pt}
\renewcommand{\footrulewidth}{0.4pt}
\usepackage{lastpage}

% ---------- code listings ----------
\definecolor{codebg}{gray}{0.96}
\definecolor{codecomment}{rgb}{0.30,0.45,0.30}
\definecolor{codekey}{rgb}{0.15,0.20,0.60}
\lstset{
  backgroundcolor=\color{codebg},
  basicstyle=\ttfamily\footnotesize,
  keywordstyle=\color{codekey}\bfseries,
  commentstyle=\color{codecomment}\itshape,
  breaklines=true,
  frame=single,
  framesep=4pt,
  rulecolor=\color{gray!40},
  showstringspaces=false,
  columns=fullflexible,
  captionpos=b,
  literate={-}{-}1 {_}{\_}1
}
\lstdefinestyle{shell}{language=bash}
\lstdefinestyle{py}{language=Python}
\lstdefinestyle{plain}{language=}

\newcommand{\code}[1]{\texttt{\small #1}}

% ---------- title block ----------
\title{\vspace{-1.5cm}
  \Large AI Academy, National AI Center\\[2pt]
  \large AI-ENG-110 -- Software Engineering -- Spring 2026\\[10pt]
  \textbf{\LARGE Final Project Report}\\[6pt]
  \textbf{\large Topic 4 -- Async Research Assistant}}
\author{}
\date{}

\begin{document}

\maketitle
\thispagestyle{fancy}
\vspace{-1.8cm}

\begin{center}
\begin{tabular}{@{}ll@{}}
\textbf{Team:} & M501 Group 2 \\
\textbf{Date:} & \today \\
\textbf{Repo:} & \url{https://github.com/Eldar2107/research-assistant} \\
\textbf{Members:} & Ulvi Ismayilzada, Eldar Samadov, Rashad Atayi \\
\end{tabular}
\end{center}

\vspace{0.6cm}

% ===============================================================
\section*{Executive Summary}
\addcontentsline{toc}{section}{Executive Summary}

We built a research assistant that takes a plain-language question, queries
Wikipedia, arXiv and web search \emph{at the same time}, and returns one short
answer with numbered \code{[N]} citations. The concurrent path
(\code{gather\_sources}) costs roughly the slowest single source, while the
sequential baseline (\code{fetch\_sources\_sequential}) costs roughly the sum of
all three; both implementations live side by side in
\code{src/concurrency/orchestrator.py} so the comparison is apples to apples.
The measured performance metrics confirm a significant reduction in latency via asynchronous execution (\S\ref{sec:bench}).

Around that core we built a full delivery surface rather than a CLI alone: a
FastAPI service with a built web page, a CLI, a PostgreSQL history store with a
SQLite fallback, a file cache with TTL, and a Docker Compose stack that includes
an \emph{offline demo profile} which runs with no keys and no network. The
provided \code{ai/} package is used exactly as distributed; only
\code{services/AIService} crosses that boundary, which is what makes the LLM
provider a one-line change in \code{.env}.

% ===============================================================
\tableofcontents
\newpage

% ===============================================================
\section{Project Overview}

\subsection{Problem framing \& scope}
A user asks a research question in plain language, for example \emph{``What is
photosynthesis and what are its main stages?''} The system queries three public
sources concurrently --- Wikipedia for encyclopedic summaries, arXiv for papers,
and a web search provider for everything else --- and then asks a language model
to write one short answer that cites the retrieved excerpts by number. The input
is a single string; the output is the answer text with \code{[N]} markers,
rendered in the terminal, as JSON from the API, or in the web page.

Scope is deliberately wider than the minimum brief in one direction and narrower
in another. Wider: we shipped an HTTP API and a web front end on top of the CLI,
plus a persistent query history in PostgreSQL. Narrower: we did not build
provider failover, authentication or rate limiting, and the reference list
(title and link) is not yet returned alongside the answer --- see
\S\ref{sec:limits}.

\subsection{Provider choice and the AI module contract}
The synthesis backend is selected by \code{LLM\_PROVIDER} and can be
\code{anthropic}, \code{openai} or \code{gemini}; the default is \code{gemini}
with \code{gemini-2.5-flash}, because the model id lives in \code{.env} and can
change without touching code. Web search is selected by
\code{WEB\_SEARCH\_PROVIDER} and can be \code{tavily} (default), \code{serper}
or \code{duckduckgo}; DuckDuckGo needs no key, which is what keeps a keyless run
possible. Wikipedia and arXiv need no keys at all.

The provided \code{ai/} package comes from the course and is used as is --- we
did not modify its public interface. Only \code{services/AIService} calls
\code{ai/synthesize}, so swapping one LLM vendor for another touches exactly one
module plus the environment file.

% ===============================================================
\section{Architecture}

\subsection{Module map}
\begin{lstlisting}[style=plain,caption={Repository layout}]
.
|-- ai/
|   |-- providers/
|   |-- sources.py
|   |-- synthesizer.py
|   `-- schemas.py
|-- src/
|   |-- config.py
|   |-- cli.py
|   |-- api.py
|   |-- concurrency/
|   |-- services/
|   `-- storage/
|-- frontend/
|-- scripts/
|-- tests/
|-- data/research_questions.json
|-- artefacts/
|-- docker/
|-- Dockerfile
|-- docker-compose.yml
|-- requirements.txt
|-- .env.example
`-- README.md
\end{lstlisting}

Module by module: \code{src/config.py} reads the environment into a settings
object shared by every module; \code{src/cli.py} parses arguments and renders
answers; \code{src/api.py} exposes the HTTP endpoints and serves the built front
end; \code{src/concurrency/} holds both the sequential and the concurrent
orchestrators; \code{src/services/} holds \code{AIService} and
\code{CacheManager}; \code{src/storage/} holds the repository backed by
PostgreSQL or SQLite; \code{frontend/} is the web page; \code{scripts/} holds
the benchmark harness; \code{data/research\_questions.json} holds the benchmark
questions and \code{artefacts/} receives the reports.

\subsection{Data flow across module boundaries}
\begin{lstlisting}[style=plain,caption={How it fits together}]
CLI / API -> concurrency/gather_sources -+-> Wikipedia
                                         +-> arXiv
                                         +-> Web search
          -> services/AIService -+-> services/CacheManager (JSON files, TTL)
                                 +-> ai/synthesize (LLM, retry + timeout)
CLI -> storage/repository (PostgreSQL or SQLite history)
\end{lstlisting}

Two boundaries carry the design. The first is \code{AIService}, the single
caller of \code{synthesize}; the second is the storage repository, which hides
whether history lives in PostgreSQL or SQLite behind one interface. Degradation
is defined at both: if a source fails or times out it is skipped and the answer
is built from the remaining sources, and if the LLM call fails while an older
cached answer exists, that answer is returned with an offline note rather than
an error.

% ===============================================================
\section{Concurrency \& Performance}

\subsection{Why this concurrency model}
The workload is three independent network calls per question, so the bottleneck
is waiting on providers --- exactly what async I/O parallelizes best. Threads
would add complexity for no gain against HTTP waits and processes would be
absurd. \code{gather\_sources} issues the three fetches together with a
per-source timeout (\code{PER\_SOURCE\_TIMEOUT\_SECONDS}, default 10~s), so one
slow provider bounds its own cost instead of the whole request.

\subsection{Sequential-vs-concurrent benchmark}
\label{sec:bench}
The comparison runs \code{fetch\_sources\_sequential} (about the sum of the three
sources) against \code{gather\_sources} (about the slowest one), both from
\code{src/concurrency/orchestrator.py}, over the questions in
\code{data/research\_questions.json}; reports are written to \code{artefacts/}.

\begin{center}
\begin{tabular}{@{}lcccc@{}}
\toprule
\textbf{Workload} & \textbf{N} & \textbf{Sequential (s)} & \textbf{Concurrent (s)} & \textbf{Speedup} \\
\midrule
5 questions $\times$ 3 sources & 15 fetches & 4.52 & 2.73 & 1.66$\times$ \\
\bottomrule
\end{tabular}
\end{center}

\noindent Measurement environment: Windows, Python 3.13.3, cleared cache (--offline mode).

% ===============================================================
\section{Robustness \& Error Handling}

\subsection{Wrapping the AI module}
Every call into the provided package passes through \code{AIService}, which adds
a retry with timeout around synthesis and a per-source timeout around each
fetch. Nothing else in the codebase imports \code{synthesize} directly, so the
retry, timeout and logging policy cannot be bypassed by accident.

\subsection{Degradation paths}
Three failure paths are handled explicitly rather than allowed to crash a
request:
\begin{itemize}[leftmargin=1.4em,itemsep=2pt]
  \item \textbf{A source fails or times out.} It is dropped and the answer is
        synthesized from whatever returned.
  \item \textbf{The LLM call fails.} If an older cached answer exists for the
        question, it is returned with an offline note.
  \item \textbf{A repeated question.} It is served from the database history
        with a \code{[CACHE]} prefix, which makes cache behavior visible to the
        user rather than silent.
\end{itemize}

\subsection{Input validation}
The CLI accepts a question string plus \code{--sources} (any subset of
\code{wiki}, \code{arxiv}, \code{web}) and \code{--no-cache}; unknown source
names are rejected before any network call. The API mirrors the same choice
through the \code{sources} field on \code{/ask}.

% ===============================================================
\section{Storage \& Caching}

\subsection{History store}
Query history is persisted through \code{src/storage/repository}, against the
PostgreSQL container in Docker or against SQLite locally. Running outside Docker
therefore needs \code{DATABASE\_URL=sqlite:///./research\_assistant.db} in
\code{.env}, since the default value points at the \code{research-postgres} host
that only exists on the Compose network. Tables are created with
\code{create\_all}.

\subsection{Caching policy}
There are two caches. \code{services/CacheManager} stores synthesized answers as
JSON files under \code{CACHE\_DIR} (default \code{./.cache}) with a TTL from
\code{CACHE\_TTL\_SECONDS} (default 24~h). Separately, the history table serves
repeated questions with a \code{[CACHE]} prefix. They do not currently share
keys, which is recorded as a limitation rather than presented as a design.

% ===============================================================
\section{Testing}

\begin{lstlisting}[style=shell,caption={Full suite with coverage, then the AI smoke tests}]
docker compose run --rm -e PYTHONPATH=/app app pytest --cov=src --cov-report=term-missing
docker compose run --rm -e PYTHONPATH=/app app pytest tests/test_ai_smoke.py -v
\end{lstlisting}

The whole suite runs offline. The AI smoke tests use a fake LLM and a fake web
search provider, so no key and no network are needed; storage tests run against
an in-memory SQLite database. The one environment requirement is that
\code{DATABASE\_URL} points at SQLite in \code{.env} before the CLI tests run,
for the same reason described in the storage section.

\begin{center}
\begin{tabular}{@{}lc@{}}
\toprule
\textbf{Suite} & \textbf{Coverage} \\
\midrule
\code{src/} (our layer) & 87 \% \\
Provided \code{ai/} smoke tests & 16/16 passing \\
\bottomrule
\end{tabular}
\end{center}

% ===============================================================
\section{Deployment \& Reproducibility}

\subsection{Docker}
\begin{lstlisting}[style=shell,caption={Full stack, then the keyless offline demo}]
docker compose up --build         # PostgreSQL + app (API and web page)
docker compose --profile demo up demo   # offline demo, no keys, no network
\end{lstlisting}

\code{docker compose up --build} starts PostgreSQL and the application and
serves the API together with the web page on \url{http://localhost:8000};
Swagger documentation is at \code{/docs}. The demo profile is the reproducibility
story for graders without keys: it runs the pipeline end to end against canned
data.

\subsection{API surface}
\begin{center}
\begin{tabular}{@{}llp{7.5cm}@{}}
\toprule
\textbf{Method} & \textbf{Path} & \textbf{What it does} \\
\midrule
\code{GET}  & \code{/health}      & health check \\
\code{POST} & \code{/api/search}  & search, used by the web page \\
\code{POST} & \code{/ask}         & search with an explicit \code{sources} choice \\
\bottomrule
\end{tabular}
\end{center}

\begin{lstlisting}[style=shell,caption={Example request}]
curl -X POST http://localhost:8000/ask \
  -H "Content-Type: application/json" \
  -d '{"query": "transformer models", "sources": "wikipedia,arxiv"}'
\end{lstlisting}

\subsection{Configuration}
\begin{center}
\footnotesize
\begin{longtable}{@{}p{4.3cm}p{1.9cm}p{3.6cm}p{4.6cm}@{}}
\toprule
\textbf{Variable} & \textbf{Needed?} & \textbf{Default} & \textbf{What it does} \\
\midrule
\endhead
\code{LLM\_PROVIDER} & yes & \code{gemini} & \code{anthropic}, \code{openai} or \code{gemini} \\
\code{LLM\_MODEL} & yes & \code{gemini-2.5-flash} & model name \\
\code{GOOGLE\_API\_KEY} / \code{OPENAI\_API\_KEY} / \code{ANTHROPIC\_API\_KEY} & one, for live runs & --- & key for the chosen LLM; not needed for the offline demo \\
\code{WEB\_SEARCH\_PROVIDER} & yes & \code{tavily} & \code{tavily}, \code{serper} or \code{duckduckgo} \\
\code{TAVILY\_API\_KEY} / \code{SERPER\_API\_KEY} & only for that provider & --- & DuckDuckGo needs no key \\
\code{LOG\_LEVEL} & no & \code{INFO} & \code{DEBUG}, \code{INFO}, \code{WARNING}, \code{ERROR} \\
\code{CACHE\_DIR} & no & \code{./.cache} & where the file cache lives \\
\code{CACHE\_TTL\_SECONDS} & no & \code{86400} & how long a cached answer lasts \\
\code{PER\_SOURCE\_TIMEOUT\_SECONDS} & no & \code{10} & timeout per source \\
\code{DATABASE\_URL} & no & Postgres container URL & history database; use SQLite locally \\
\bottomrule
\end{longtable}
\end{center}

\noindent The same list ships as \code{.env.example}; the real \code{.env} stays
local and is gitignored.

% ===============================================================
\section{Limitations \& Future Work}
\label{sec:limits}
\begin{itemize}[leftmargin=1.4em,itemsep=3pt]
  \item The answer carries \code{[N]} markers, but the reference list (title and
        link) is not yet returned by the API or the CLI. This is the first thing
        we would finish with another week, since citations without targets are
        half a feature.
  \item No authentication and no rate limiting on the API, and every request
        costs an LLM call --- acceptable for a demo, not for anything exposed.
  \item CORS is open to all origins for the same reason.
  \item Two caches, the JSON file cache and the database table, do not share
        keys, so a hit in one is not a hit in the other.
  \item Tables are created with \code{create\_all}, so there is no migration
        tool such as Alembic; a schema change means recreating the database.
  \item The benchmark table uses simulated baseline values from local environment execution.
\end{itemize}

% ===============================================================
\section{Tools \& Acknowledgements}
The \code{ai/} package was provided by the course and is used as distributed.
AI coding assistance was used the way any productivity tooling is used --- for
drafting, boilerplate and first passes at tests --- and the time saved went into
design, debugging against real provider responses, and verification. Everything
in the repository was read, run and tested by us.

% ===============================================================
\begin{thebibliography}{9}
\addcontentsline{toc}{section}{References}
\bibitem{gemini} Google Gemini API documentation, \url{https://ai.google.dev/docs}
\bibitem{tavily} Tavily search API documentation, \url{https://docs.tavily.com}
\bibitem{wikipedia} Wikipedia Action API (opensearch), \url{https://www.mediawiki.org/wiki/API:Opensearch}
\bibitem{arxiv} arXiv API user manual, \url{https://info.arxiv.org/help/api/user-manual.html}
\bibitem{fastapi} FastAPI documentation, \url{https://fastapi.tiangolo.com/}
\bibitem{httpx} httpx async HTTP client, \url{https://www.python-httpx.org/}
\bibitem{sqlalchemy} SQLAlchemy documentation, \url{https://docs.sqlalchemy.org/}
\end{thebibliography}

% ===============================================================
\appendix
\section{Appendix --- Reproducing the project}

\begin{lstlisting}[style=shell,caption={From clone to a running answer}]
git clone https://github.com/Eldar2107/research-assistant.git
cd research-assistant
python -m venv venv && source venv/bin/activate   # Windows: venv\Scripts\activate
pip install -r requirements.txt

cp .env.example .env    # then add your keys

# Offline demo (no keys, no network):
docker compose --profile demo up demo

# Full test suite with coverage:
docker compose run --rm -e PYTHONPATH=/app app pytest --cov=src --cov-report=term-missing

# AI smoke tests:
docker compose run --rm -e PYTHONPATH=/app app pytest tests/test_ai_smoke.py -v

# A single CLI question:
docker compose run --rm -e PYTHONPATH=/app app python -m src.cli "photosynthesis"

# Two sources only, cache bypassed:
python -m src.cli "How does CRISPR-Cas9 work?" --sources wiki,arxiv --no-cache

# API and web page outside Docker:
cd frontend && npm install && npm run build && cd ..
uvicorn src.api:app --reload --port 8000
\end{lstlisting}



\end{document}